\chapter{La carte des fournisseurs}\label{ch:nw:the-vendor-map}

La liste des \omterm{def:nw:the-vendor-map:isv}{éditeurs de logiciels indépendants} d'Eurex compte trente-trois sociétés, chacune repérée par la partie d'une société de trading qu'elle sert~: back office, middle office, front office.
Ses autres listes recensent les fournisseurs de données de marché, les courtiers et les fournisseurs d'accès
sponsorisé. Une nouvelle société de trading achète l'essentiel de sa pile technique, et sa première carte du marché est une carte des fournisseurs~: auprès de qui elle achète
sa \omterm{def:nw:colocation-products-and-how-they-are-sold:colo}{colocation}, quel \omterm{def:nw:the-vendor-map:tickerplant}{ticker plant} décode ses flux, dont les serveurs font tourner ses stratégies, dont les horloges horodatent ses enregistrements.

Les chapitres précédents ont rencontré ces fournisseurs un produit à la fois~: commutateurs et dispositifs de couche~1 au chapitre~2, \omterm{def:nw:networking-for-trading:frame}{cartes réseau}
au chapitre~3, horloges au chapitre~4, capture au chapitre~5, cartes FPGA aux chapitres~6 et~7, serveurs au chapitre~8, \omterm{def:nw:colocation-products-and-how-they-are-sold:colo}{colocation} et
routes aux chapitres~9 à~15. Ce chapitre les place sur une seule carte, par catégorie, rattache la pile d'une société à cette carte sous forme de graphe, et
mesure ce que la carte montre~: à quel point les dépenses de la société sont concentrées et quelles défaillances de fournisseurs l'empêcheraient de trader. Les fournisseurs
ne sont nommés qu'à partir de la liste d'une plateforme ou de leurs propres publications~; la pile de l'exemple est fictive.

\section{Données de marché}

\begin{definition}[Ticker plant]\label{def:nw:the-vendor-map:tickerplant}
Un \emph{ticker plant}\index{ticker plant} est un système, acheté ou construit, qui reçoit une fois les flux de données de marché d'une société, les décode
et les normalise dans un format unique, tient à jour les carnets d'ordres, et distribue les résultats aux applications de trading, de risque et
de suivi de la société par son réseau interne.
\end{definition}

Une société peut décoder chaque flux à l'intérieur de chaque stratégie (le feed handler du livre~13) ou une seule fois, de manière centralisée, dans un \omterm{def:nw:the-vendor-map:tickerplant}{ticker plant}.
Les fournisseurs vendent les deux. Exegy décrit les \omterm{def:nw:the-vendor-map:tickerplant}{ticker plants} centralisés comme «~la réponse architecturale dominante aux limites de passage à l'échelle des
feed handlers embarqués~», et en nomme le coût~: «~la complexité de distribution et la surcharge de latence~»
(\cref{dat:nw:the-vendor-map:lists}). La page produit de Pico définit le \omterm{def:nw:the-vendor-map:tickerplant}{ticker plant} comme un logiciel qui se connecte aux flux,
normalise les messages et tient les carnets en mémoire, pour plus de 230 plateformes. Le choix est la question du multicast du chapitre~1
transformée en achat~: un seul décodage et un second réseau, ou plusieurs décodages et aucun.

\section{Logiciels d'exécution et gestion des ordres}

\begin{definition}[Éditeur de logiciels indépendant]\label{def:nw:the-vendor-map:isv}
Un \emph{éditeur de logiciels indépendant}\index{editeur de logiciels independant@éditeur de logiciels indépendant} (independent software vendor) est une société qui vend des logiciels de trading, de gestion des ordres,
de risque ou de post-marché à des sociétés de trading, et qu'une bourse certifie pour se connecter à ses systèmes en leur nom, sans être
elle-même membre.
\end{definition}

Les bourses publient leurs éditeurs certifiés. La liste d'Eurex indique les services de chaque société par bureau~: Trading Technologies pour le back
et le front office, Broadridge et OnixS pour le front office, parmi beaucoup d'autres. Pour une nouvelle société, la liste est à la fois une liste d'achats et
une contrainte~: un fournisseur qui n'est pas certifié pour une plateforme ne peut pas y connecter la société, et la certification de chaque plateforme est un
projet distinct (livre~13).

\begin{dated}{2026-09}{Fournisseurs sur la carte, d'après les listes des plateformes et les publications des fournisseurs}\label{dat:nw:the-vendor-map:lists}
\begin{itemize}
\item \omterm{def:nw:the-vendor-map:isv}{Éditeurs de logiciels indépendants} d'Eurex~: 33 sociétés, chacune avec ses bureaux servis, parmi lesquelles Trading Technologies (back
  et front office), Broadridge (front office), OnixS (front office).
\item Données de marché~: Exegy (\omterm{def:nw:the-vendor-map:tickerplant}{ticker plants} centralisés, décodage des flux sur FPGA)~; Pico (InRush Ticker Plant, plus de 230 plateformes).
\item Synchronisation~: Meinberg (\omterm{def:nw:time-synchronisation:gm}{horloges grand maître} IEEE 1588, serveurs de temps LANTIME). Suivi~: Pico (Corvil Analytics).
\item Issus des chapitres précédents~: commutateurs (Cisco, Arista), \omterm{def:nw:networking-for-trading:frame}{cartes réseau} et cartes FPGA (AMD), serveurs overclockés (Blackcore,
  Business Systems International), \omterm{def:nw:colocation-products-and-how-they-are-sold:colo}{colocation} (Equinix), une route hertzienne (Quincy Data).
\end{itemize}
\end{dated}

\section{Matériel réseau et serveurs}

Les catégories de matériel sont celles mesurées aux chapitres~2 à~8, et chacune compte une poignée de fournisseurs que les sociétés à faible latence utilisent~: un
fournisseur de commutateurs pour la couche d'agrégation et un fournisseur de dispositifs de couche~1 pour la duplication (fan-out)~; une \omterm{def:nw:networking-for-trading:frame}{carte réseau} qui prend en charge le kernel bypass~; une
carte FPGA pour le chemin matériel~; des serveurs overclockés. Une société standardise en général sur un seul fournisseur par catégorie, pour les pièces de rechange, les
compétences et les pilotes~; ce choix fait de ce fournisseur un point de défaillance unique, que le plan de reprise du chapitre~28 doit couvrir.

\section{Synchronisation et suivi}

Synchronisation et suivi sont les catégories qu'une société peut survivre à perdre pendant une heure et ne doit pas perdre pendant une journée. Une \omterm{def:nw:time-synchronisation:gm}{horloge grand maître}
alimente chaque horodatage que les régulateurs demandent (chapitre~4)~; l'analyse réseau et les \omterm{def:nw:capturing-and-monitoring-the-network:capture}{boîtiers de capture} (chapitre~5) voient ce que les propres systèmes
de la société ne voient pas. Ni l'un ni l'autre ne se trouve sur le chemin qui mène des données aux ordres, donc aucun des deux n'arrête le trading~; tous deux empêchent
la société de prouver ce qu'elle a fait, ce qui constitue une autre forme d'arrêt.

\section{Lire la carte~: enfermement, concentration et sortie}

Une pile est un graphe~: les données de marché y entrent, traversent des composants et en sortent sous forme d'ordres. Chaque composant est fourni par un ou
plusieurs fournisseurs, ou construit en interne. Deux questions en découlent. Quels composants, en défaillant seuls, coupent tous les chemins allant des données aux ordres~? Et
quels fournisseurs, en défaillant, emportent chaque composant qu'ils sont seuls à fournir et le coupent ainsi~?

\omcode{code/firm/vendormap/firm_vendormap.py}{60}{83}{Chemins des données aux ordres, composants qui les coupent seuls, et fournisseurs qui le font.}\label{lst:nw:the-vendor-map:graph}

\begin{proposition}[Points de défaillance des fournisseurs]\label{prop:nw:the-vendor-map:vendors}
Un fournisseur interrompt le trading si et seulement si les composants qu'il est seul à fournir forment une coupure entre les données et les ordres~; un composant fourni par
deux fournisseurs n'est jamais retiré par la défaillance d'un seul d'entre eux, de sorte qu'une double source pour un composant le retire de la coupure de chaque fournisseur, mais pas de
celle de son propre fournisseur.
\end{proposition}

\begin{proof}
La défaillance d'un fournisseur retire exactement les composants qui n'ont pas d'autre fournisseur~; le trading s'arrête si et seulement si aucun chemin allant des données aux
ordres ne survit à leur retrait. Un composant à double source survit à la défaillance de l'un ou l'autre fournisseur, mais le composant lui-même, en défaillant pour
une autre cause (son logiciel, sa configuration), coupe encore le chemin s'il est la seule route.
\end{proof}

La société de l'exemple de ce chapitre opère depuis Secaucus. Sa pile~: \omterm{def:nw:colocation-products-and-how-they-are-sold:colo}{colocation} et \omterm{def:nw:colocation-products-and-how-they-are-sold:mmr}{interconnexions} fournies par Equinix~; une duplication de couche~1 (fan-out) chez Arista
et des commutateurs chez Cisco, en parallèle~; un \omterm{def:nw:the-vendor-map:tickerplant}{ticker plant} Exegy et des feed handlers internes, en parallèle~; des serveurs Blackcore avec des \omterm{def:nw:networking-for-trading:frame}{cartes réseau}
AMD~; sa propre passerelle d'exécution~; et, hors du chemin, la gestion des ordres de Trading Technologies, les \omterm{def:nw:time-synchronisation:gm}{horloges grand maître} de Meinberg, l'analyse réseau
de Pico et une route hertzienne de Quincy Data. Elle y consacre 3.19~millions de dollars par an (chiffres fictifs).

\begin{omfigure}
\begin{tikzpicture}[font=\footnotesize,node distance=4mm,x=1.13cm,
  spof/.style={draw=omThm,fill=omThm!25,rounded corners,align=center,minimum height=8mm,inner sep=2pt,text width=2.3cm},
  ok/.style={draw=omDef,fill=omDef!10,rounded corners,align=center,minimum height=8mm,inner sep=2pt,text width=2.3cm},
  off/.style={draw=gray,fill=gray!8,rounded corners,align=center,minimum height=8mm,inner sep=2pt,text width=2.3cm,dashed}]
\node[spof] (colo) at (0,0) {colocation\\Equinix};
\node[spof] (xc) at (2.6,0) {interconnexions\\Equinix};
\node[ok] (l1) at (5.2,0.75) {duplication couche 1\\Arista};
\node[ok] (sw) at (5.2,-0.75) {commutateurs\\Cisco};
\node[ok] (tp) at (7.8,0.75) {ticker plant\\Exegy};
\node[ok] (fh) at (7.8,-0.75) {feed handlers\\interne};
\node[spof] (srv) at (10.4,0) {serveurs\\Blackcore};
\node[spof] (nic) at (10.4,-2.0) {cartes réseau\\AMD};
\node[spof] (gw) at (7.8,-2.6) {passerelle\\interne};
\node[off] (oms) at (0,-2.2) {ordres (OMS)\\Trading Technologies};
\node[off] (clk) at (2.6,-2.2) {horloges grand maître\\Meinberg};
\node[off] (mon) at (5.2,-2.6) {analyse\\Pico};
\draw[->] (colo) -- (xc); \draw[->] (xc) -- (l1); \draw[->] (xc) -- (sw);
\draw[->] (l1) -- (tp); \draw[->] (sw) -- (fh); \draw[->] (l1) -- (fh); \draw[->] (sw) -- (tp);
\draw[->] (tp) -- (srv); \draw[->] (fh) -- (srv); \draw[->] (srv) -- (nic); \draw[->] (nic) -- (gw);
\end{tikzpicture}

\omcaption{La pile de l'entreprise exemple, des données (à gauche) aux ordres (la passerelle). En rouge plein~: les composants dont la seule défaillance interrompt
la négociation~; en bleu~: les composants disposant d'une alternative parallèle~; en gris pointillé~: hors du chemin de négociation. Pile synthétique~; éditeurs comme dans
l'\cref{dat:nw:the-vendor-map:lists}. Données~: \texttt{nw\_vendors.report()}.}\label{fig:nw:the-vendor-map:stack}
\end{omfigure}

Cinq composants coupent seuls le chemin~: la \omterm{def:nw:colocation-products-and-how-they-are-sold:colo}{colocation}, les \omterm{def:nw:colocation-products-and-how-they-are-sold:mmr}{interconnexions}, les serveurs de stratégie, leurs \omterm{def:nw:networking-for-trading:frame}{cartes réseau} et la passerelle interne.
Trois éditeurs le font~: AMD, Blackcore et Equinix~; Arista, Cisco et Exegy ne le font pas, car chacun dispose d'un chemin parallèle. L'indice de Herfindahl--Hirschman
des dépenses de l'entreprise par éditeur (parts de la dépense totale, travaux internes exclus de l'indice) vaut 1\,149 sur une échelle de zéro à
10\,000, ce qui se lit comme une dispersion~; mais l'indice mesure les dépenses, non la dépendance~: le
premier éditeur de l'entreprise en dépenses, Equinix avec 22.6~\%, est un point de défaillance, alors que le second, Quincy Data avec 15.0~\%, ne l'est pas.

\begin{omfigure}
\begin{tikzpicture}
\begin{axis}[xbar,width=10.5cm,height=6.4cm,bar width=6pt,xmin=0,xmax=27,ytick=data,y dir=reverse,
  yticklabels={Equinix,Quincy Data,en interne,Exegy,Trading Technologies,Blackcore,Pico,Cisco,Arista,AMD,Meinberg},
  xlabel={part des dépenses annuelles (\%)},nodes near coords,every node near coord/.append style={font=\scriptsize},
  nodes near coords={\pgfmathprintnumber[fixed,fixed zerofill,precision=1]{\pgfplotspointmeta}},
  tick label style={font=\footnotesize},label style={font=\small},grid=major,grid style={gray!25},table/col sep=comma]
\addplot[fill=omDef!55,draw=omDef] table[y=x,x=share]{figdata/networks/27-the-vendor-map/shares.csv};
\end{axis}
\end{tikzpicture}

\omcaption{Les dépenses annuelles de l'entreprise exemple par éditeur (chiffres synthétiques)~: la dépense est répartie, la dépendance ne l'est pas. Données~:
\texttt{fig\_vendors.py}, \texttt{nw\_vendors.report()}.}\label{fig:nw:the-vendor-map:shares}
\end{omfigure}

Ajouter un second éditeur de serveurs sur la moitié des serveurs ramène l'indice à 1\,119 et retire Blackcore des points éditeurs~; les serveurs de
stratégie restent un composant unique qui peut tomber en panne pour d'autres raisons. Les délais de sortie des cinq composants à point unique s'additionnent pour atteindre 48 semaines~:
quitter Equinix seul représente un semestre. C'est le sens pratique de l'enfermement auprès d'un éditeur (livre~16, chapitre~20, le traite comme une
décision commerciale)~: non pas une clause de contrat, mais les semaines nécessaires pour remplacer un composant du chemin de négociation.

\begin{method}[Lire la carte des éditeurs d'une entreprise]\label{met:nw:the-vendor-map:read}
\begin{enumerate}
\item Lister chaque composant du chemin allant des données de marché aux ordres, et chaque composant hors de ce chemin~; rattacher chacun à ses éditeurs, avec la source
  qui nomme l'éditeur et la date.
\item Dessiner le graphe~; repérer les composants uniques et les éditeurs dont la défaillance le coupe.
\item Calculer la concentration des dépenses, et la lire à côté des points éditeurs~: elles répondent à des questions différentes.
\item Pour chaque point éditeur, estimer le délai de sortie et décider~: double source, pièces de rechange, ou accepter et couvrir dans le plan de reprise
  (chapitre~28).
\item Exporter la nomenclature vers le plan de connectivité (chapitre~29) et relancer la carte à chaque renouvellement de contrat.
\end{enumerate}
\end{method}

\section{Tutoriel~: un éditeur de trop}\label{tut:nw:the-vendor-map:tut}

\begin{tutorial}
\textbf{Objectif.} Construire la carte des éditeurs sous forme de données, rattacher une pile et repérer ses points de défaillance uniques.
\textbf{État final~:} \cref{fig:nw:the-vendor-map:stack,fig:nw:the-vendor-map:shares}.
\begin{tutsteps}
\item \textbf{Inventaire.} \texttt{firm\_vendormap.load\_inventory} lit les lignes datées des éditeurs, chacune avec sa source dans le registre.
\item \textbf{Pile.} \texttt{nw\_vendors.COMPONENTS} et \texttt{EDGES} décrivent l'entreprise exemple~; \texttt{Stack} les contient.
\item \textbf{Mesures.} \texttt{single\_points}, \texttt{vendor\_points} (le \cref{lst:nw:the-vendor-map:graph}), \texttt{hhi} et
  \texttt{bom}~; \texttt{second\_server\_vendor} pour le changement.
\end{tutsteps}
\textbf{Pour aller plus loin.} Pondérer chaque point éditeur par son délai de sortie et la probabilité de sa défaillance, puis classer les corrections par
risque retiré par dollar.
\end{tutorial}

\section{Construction~: la carte des éditeurs}\label{bld:nw:the-vendor-map:vendormap}

\begin{build}
\bfield{Rôle} Un inventaire des éditeurs et la pile d'une entreprise sous forme de graphe, avec concentration, points de défaillance uniques et une
nomenclature~: les lignes éditeurs du plan du chapitre~29.
\bfield{Interface} \texttt{firm\_vendormap}~: \texttt{InventoryRow}, \texttt{load\_inventory}, \texttt{Component}, \texttt{Stack}
(\texttt{connected}, \texttt{single\_points}, \texttt{vendor\_points}, \texttt{shares}, \texttt{hhi}, \texttt{bom}).
\bfield{Règles} Chaque ligne éditeur nomme sa source dans le registre et sa date~; la pile et ses dépenses appartiennent à l'entreprise (ici synthétiques).
\bfield{Tests d'acceptation} \texttt{code/firm/vendormap/tests/}~: sources de l'inventaire, un graphe à trois composants vérifié à la main, et
composants internes exclus des points éditeurs.
\bfield{Pour aller plus loin} Probabilités de défaillance et coûts de sortie~; sous-fournisseurs partagés (deux éditeurs sur une même puce).
\end{build}

\begin{omsources}
\item Eurex, listes de participants~: éditeurs de logiciels indépendants.
\item Exegy, \emph{Design Patterns for Market Data, Part~2: Centralized Ticker Plant}~; Pico, InRush Ticker Plant~; Meinberg,
  gamme de produits.
\item Les chapitres~2 à~14 de ce livre, pour les éditeurs de matériel, de colocation et de routes.
\end{omsources}

\section{Exercices}

\begin{exercise}[$\star$]\label{exo:nw:the-vendor-map:1}
Que fait un \omterm{def:nw:the-vendor-map:tickerplant}{ticker plant}, et que coûte sa centralisation~?
\end{exercise}

\begin{exercise}[$\star$]\label{exo:nw:the-vendor-map:2}
Qu'est-ce qu'un \omterm{def:nw:the-vendor-map:isv}{éditeur de logiciels indépendant}, et pourquoi une bourse le certifie-t-elle~?
\end{exercise}

\begin{exercise}[$\star$]\label{exo:nw:the-vendor-map:3}
Pourquoi les \omterm{def:nw:time-synchronisation:gm}{horloges grand maître} n'apparaissent-elles pas parmi les points de défaillance de l'entreprise exemple, et pourquoi l'entreprise doit-elle quand même s'en soucier~?
\end{exercise}

\begin{exercise}[$\star\star$]\label{exo:nw:the-vendor-map:4}
En utilisant la \cref{prop:nw:the-vendor-map:vendors}, pourquoi Exegy n'est-il pas un point de défaillance éditeur pour l'entreprise exemple~?
\end{exercise}

\begin{exercise}[$\star\star$]\label{exo:nw:the-vendor-map:5}
Calculer l'indice de l'entreprise exemple avec les dépenses d'Equinix doublées.
\end{exercise}

\begin{exercise}[$\star\star$]\label{exo:nw:the-vendor-map:6}
Le second éditeur de serveurs de l'entreprise achète ses \omterm{def:nw:networking-for-trading:frame}{cartes réseau} au même fabricant. Que manque-t-il à la carte~?
\end{exercise}

\begin{exercise}[$\star\star\star$]\label{exo:nw:the-vendor-map:7}
\emph{Programmation.} Avec \texttt{firm\_vendormap.Stack}, ajouter un second site de \omterm{def:nw:colocation-products-and-how-they-are-sold:colo}{colocation}, chez un autre opérateur, avec ses propres \omterm{def:nw:colocation-products-and-how-they-are-sold:mmr}{interconnexions} alimentant les mêmes serveurs.
Quels points éditeurs subsistent~?
\end{exercise}

\begin{exercise}[$\star\star\star$]\label{exo:nw:the-vendor-map:8}
\emph{Trouver la faille.} «~Notre premier éditeur ne représente que 23~\% de nos dépenses, donc nous ne dépendons de personne.~»
\end{exercise}

\section{Problème~: Un éditeur de trop}

\begin{problem}[{Problème du week-end --- concentration et points de défaillance uniques chez les éditeurs}]\label{pb:nw:the-vendor-map:1}
Une nouvelle société de trading opère depuis Secaucus avec la pile de l'exemple du chapitre. Son conseil d'administration demande quels éditeurs pourraient l'empêcher de trader et ce qu'il faudrait
faire pour changer cela.

\medskip\noindent\textbf{Partie I --- La carte.}
\begin{enumerate}
\item Quelles catégories la société achète-t-elle, et auprès de qui~?
\item Quelles sources désignent chaque éditeur~?
\item Combien la société dépense-t-elle en un an, et auprès de qui surtout~?
\item Quel est l'indice de concentration~?
\end{enumerate}

\medskip\noindent\textbf{Partie II --- Le graphe.}
\begin{enumerate}[resume]
\item Quels composants coupent seuls le chemin~?
\item Quels éditeurs le font~?
\item Pourquoi Arista, Cisco et Exegy n'en font-ils pas partie~?
\item Quels sont les temps de sortie des composants à point unique~?
\end{enumerate}

\medskip\noindent\textbf{Partie III --- Les modifications.}
\begin{enumerate}[resume]
\item Que change un second éditeur de serveurs~?
\item Que faudrait-il pour retirer AMD des points éditeurs~?
\item Pourquoi Equinix ne peut-il pas être retiré à bon marché~?
\item Que fait le logiciel hors chemin à la société s'il tombe en panne~?
\end{enumerate}

\medskip\noindent\textbf{Partie IV --- Le verdict.}
\begin{enumerate}[resume]
\item Énoncer le \emph{résultat nommé}~: l'indice de concentration de la pile de la société et la liste des composants dont la panne de l'éditeur unique
  interrompt la cotation.
\item Quelles entrées sont publiées, et lesquelles sont propres à la société~?
\item Pourquoi l'indice et les points éditeurs ne concordent-ils pas~?
\item Que doit contenir le plan de reprise du chapitre~28~?
\item Que doit demander le conseil à chaque renouvellement de contrat~?
\item Comment la carte change-t-elle pour une société qui construit son propre \omterm{def:nw:the-vendor-map:tickerplant}{ticker plant}~?
\item Où cela s'inscrit-il dans le plan du chapitre~29~?
\item En une phrase~: qu'est-ce qui rend un éditeur dangereux~?
\end{enumerate}
\end{problem}

\section{Questions d'entretien}

\begin{interviewq}[$\star$ \iqroles{développeur}]\label{iq:nw:the-vendor-map:1}
\omterm{def:nw:the-vendor-map:tickerplant}{Ticker plant} ou feed handlers embarqués~: quels sont les compromis~?
\end{interviewq}

\begin{interviewq}[$\star\star$ \iqroles{développeur}]\label{iq:nw:the-vendor-map:2}
Comment trouveriez-vous les points de défaillance uniques dans l'infrastructure d'une société de trading~?
\end{interviewq}

\begin{interviewq}[$\star\star$ \iqroles{développeur, chercheur}]\label{iq:nw:the-vendor-map:3}
Votre éditeur de \omterm{def:nw:the-vendor-map:tickerplant}{ticker plant} annonce la fin du support de votre version dans six mois. Que faites-vous~?
\end{interviewq}

\begin{interviewq}[$\star\star$ \iqroles{développeur}]\label{iq:nw:the-vendor-map:4}
Pourquoi une société standardiserait-elle sur un seul éditeur de serveurs, et quand ne devrait-elle pas le faire~?
\end{interviewq}

\begin{interviewq}[$\star\star$ \iqroles{trader}]\label{iq:nw:the-vendor-map:5}
Votre éditeur de gestion d'ordres est en panne pendant une heure. Qu'est-ce qui fonctionne encore, et que faites-vous~?
\end{interviewq}

\begin{interviewq}[$\star\star\star$ \iqroles{développeur, chercheur}]\label{iq:nw:the-vendor-map:6}
Concevez la stratégie d'éditeurs d'une nouvelle société de trading pour ses deux premières années~: quoi acheter, quoi construire, et comment garder les sorties ouvertes.
\end{interviewq}
